\documentclass[border=5pt]{standalone}
\usepackage[american]{circuitikz}
\usetikzlibrary{calc,positioning}
\standaloneenv{circuitikz}

\begin{document}
\begin{circuitikz}[
  american,
  every node/.style={font=\scriptsize},
  wire/.style={line width=0.7pt},
  netlabel/.style={font=\ttfamily\scriptsize},
  blocklabel/.style={font=\sffamily\scriptsize, align=center},
  note/.style={font=\sffamily\tiny, align=center, text=black!65},
]
\ctikzset{
  bipoles/length=1.0cm,
  resistors/scale=0.75,
  capacitors/scale=0.70,
  sources/scale=0.92,
}

% Fully differential amplifier body, drawn explicitly to keep the TI pin order.
\coordinate (amp_tip) at (3.8,0);
\coordinate (amp_nw) at (2.25,1.25);
\coordinate (amp_sw) at (2.25,-1.25);
\draw[fill=black!3, line width=0.8pt]
  (amp_nw) -- (amp_tip) -- (amp_sw) -- cycle;
\node[font=\sffamily\small] at (2.78,0) {THS4541};

\coordinate (inp)  at (2.25,0.58);
\coordinate (inn)  at (2.25,-0.58);
% Intersections of y=+/-0.58 with the two sloped output edges.
\coordinate (outm) at (3.08,0.58);
\coordinate (outp) at (3.08,-0.58);
\node[font=\sffamily\tiny, right] at (inp) {$+$};
\node[font=\sffamily\tiny, right] at (inn) {$-$};
\node[font=\sffamily\tiny, left] at (outm) {$-$};
\node[font=\sffamily\tiny, left] at (outp) {$+$};

% 33250A single-ended sinusoidal source.  The 50-ohm output impedance is
% included in the source annotation rather than drawn as a separate resistor.
\coordinate (src_top) at (-4.35,0.58);
\draw[wire] (src_top) to[sV, l_=\texttt{$v_{\mathrm{AWG}}(t)$}]
  ++(0,-1.15) coordinate (src_gnd);
\node[ground] at (src_gnd) {};
\node[blocklabel, align=center] at (-4.35,1.43)
  {Keysight 33250A\\$R_{\mathrm{out}}=50\,\Omega$};

% Driven positive input: coax termination and 499-ohm gain resistor.
\coordinate (vinext) at (-2.65,0.58);
\draw[wire] (src_top) -- (vinext) node[circ]{};
\node[netlabel, above] at (vinext) {VIN\_P\_EXT};
\draw[wire] (vinext) to[R, l_=$54.9\,\Omega$] ++(0,-1.15)
  node[ground]{};
\draw[wire] (vinext) to[R, l=$499\,\Omega$] (inp);

% Single-ended negative reference: 50 ohms in parallel with 54.9 ohms.
\coordinate (vref) at (-1.30,-0.58);
\draw[wire] (vref) node[circ]{} to[R, l=$499\,\Omega$] (inn);
\node[netlabel, above] at (vref) {VREF $\approx 0\,\mathrm{V}$};
\draw[wire] ($(vref)+(-0.25,0)$) -- ++(0,-0.25)
  to[R, l_=$50\,\Omega$] ++(0,-0.95) node[ground]{};
\draw[wire] ($(vref)+(0.25,0)$) -- ++(0,-0.25)
  to[R, l=$54.9\,\Omega$] ++(0,-0.95) node[ground]{};
% Matched 499-ohm feedback paths.
\coordinate (fb_left) at (1.90,0);
\coordinate (fb_right) at (3.60,0);
\draw[wire] (inp) -- (1.90,0.58) -- (1.90,1.82)
  to[R, l=$499\,\Omega$] (3.60,1.82)
  -- (3.60,0.58) -- (outm);
\draw[wire] (1.90,1.82) -- (1.90,2.48)
  to[C, l=$51\,\mathrm{pF}$] (3.60,2.48)
  -- (3.60,1.82);
\draw[wire] (inn) -- (1.90,-0.58) -- (1.90,-1.82)
  to[R, l_=$499\,\Omega$] (3.60,-1.82)
  -- (3.60,-0.58) -- (outp);
\draw[wire] (1.90,-1.82) -- (1.90,-2.48)
  to[C, l_=$51\,\mathrm{pF}$] (3.60,-2.48)
  -- (3.60,-1.82);

% External common-mode control, shown on the FDA's dedicated centre input.
\draw[wire] (0.55,0) -- (2.25,0);
\node[netlabel, fill=white, inner sep=1.5pt] at (1.35,0) {VOCM\_EXT};
\node[font=\sffamily\tiny, below] at (1.35,-0.08) {$0.615\,\mathrm{V}$};

% ADC-facing differential RC network from the PCB.
\coordinate (vin_n_adc) at (6.45,0.58);
\coordinate (vin_p_adc) at (6.45,-0.58);
\draw[wire] (outm) to[R, l=$499\,\Omega$] (vin_n_adc)
  -- (7.66,0.58);
\draw[wire] (outp) to[R, l_=$499\,\Omega$] (vin_p_adc)
  -- (7.66,-0.58);
\draw[wire] (vin_n_adc) to[C, l=$51\,\mathrm{pF}$] (vin_p_adc);
\node[netlabel, above] at (7.05,0.58) {VIN\_N\_ADC};
\node[netlabel, below] at (7.05,-0.58) {VIN\_P\_ADC};

% FRIDA ADC symbol and serialized comparator output.
\coordinate (adc_tip) at (7.25,0);
\coordinate (adc_nw) at (8.10,1.20);
\coordinate (adc_ne) at (10.35,1.20);
\coordinate (adc_se) at (10.35,-1.20);
\coordinate (adc_sw) at (8.10,-1.20);
\draw[fill=black!5, line width=0.8pt]
  (adc_tip) -- (adc_nw) -- (adc_ne) -- (adc_se) -- (adc_sw) -- cycle;
\node[font=\ttfamily\bfseries\small] at (8.85,0) {frida\_adc};
\draw[wire, -{Triangle[length=1.5mm,width=1.8mm]}]
  (10.35,0) -- (11.45,0) node[netlabel, right]{comp\_out};
\node[note, below] at (10.88,-0.08) {20\,Mbaud};

% Low-frequency gain and ideal first-order bandwidth estimates.
\node[note, align=left, anchor=north west] at (6.85,-1.65) {
  $\alpha=R_F/R_G=499/499=1$\\
  $Z_A=2R_G\frac{1+\alpha}{2+\alpha}=665\,\Omega$ (THS4541 Eq. 12)\\
  $Z_{\mathrm{in}}=54.9\,\Omega\parallel Z_A=50.7\,\Omega$\\
  $A_{\mathrm{eff}}\approx\alpha\frac{Z_{\mathrm{in}}}{50\,\Omega+Z_{\mathrm{in}}}=0.504\ \mathrm{V/V}$\\[1pt]
  $H(s)\approx\frac{A_{\mathrm{eff}}}
    {(1+sR_FC_F)(1+s\,2R_OC_O)}$\\
  $f_{\mathrm{FB}}=[2\pi R_FC_F]^{-1}=6.25\,\mathrm{MHz}$\\
  $f_{\mathrm{OUT}}=[2\pi(2R_O)C_O]^{-1}=3.13\,\mathrm{MHz}$\\
  $f_{-3\mathrm{dB}}\approx2.62\,\mathrm{MHz}$ (ideal cascaded poles)};

\end{circuitikz}
\end{document}
